\chapter{Linear and Regularised Baselines}\label{ch:ml:linear-and-regularised-baselines}

A team spends a quarter on gradient-boosted trees for the cross-section of 500 stocks. Out of sample, over ten years,
their long-short decile portfolio earns a Sharpe ratio of 1.90 after costs. A lasso regression on the same forty ranked
characteristics, with its penalty chosen by cross-validation, is fitted in a second and earns 1.91. The trees forecast
better (an \omterm{def:ml:why-financial-machine-learning-is-different:r2}{out-of-sample R-squared} of 0.47\% against 0.32\%), and on a portfolio that trades the extremes of the
ranking the difference is gone. This chapter builds the model every later model must beat, and measures when and
why it is hard to beat.

\section{The model to beat}

\begin{definition}[Baseline model]\label{def:ml:linear-and-regularised-baselines:baseline}
A \emph{baseline model}\index{baseline model} is the simplest credible model of a task, fitted and evaluated with the
same data, splits, target and report as the candidate: any claim that a model works is a claim that it beats the
baseline by more than the noise of the comparison.
\end{definition}

The baseline is a discipline, not a formality. It catches the leaks of chapter~3 (a baseline that scores as well as the
candidate on a target that should be hard points at the pipeline), it prices complexity, and it is the model that runs
in production while the candidate is being validated. \texttt{firm.mlbase} is the harness every cross-sectional model of
the book goes through: features ranked by date, the target demeaned by date (chapter~1), a fit on the training months,
and one report for the test months: R-squared against zero, the mean rank information coefficient with its HAC $t$
(Book~7, chapter~6), and the Sharpe ratio of an equal-weighted long-short decile portfolio, gross and net of 20 basis
points per unit traded.

The chapter's panel is \texttt{firm.mlsynth}'s with forty characteristics, six of them carrying a planted expected return
(half of its variance nonlinear), and one addition that matters: every characteristic also carries a factor return of
zero mean and 2\% monthly volatility. A portfolio sorted on characteristics then carries factor risk, as real ones do,
and its Sharpe ratio comes down to the published range: Gu, Kelly and Xiu's equal-weighted decile portfolios reach 2.45
for their best network and 0.83 for a three-variable linear model. Without the factor returns the same models would
report Sharpe ratios near 5.

\section{Regularised regression on ranked features}

Ordinary least squares, ridge regression, the lasso and the elastic net are Book~4's (chapter~16). What this chapter
adds is how their penalty is chosen and what the choice looks like on market data.

\begin{definition}[Regularisation path]\label{def:ml:linear-and-regularised-baselines:path}
The \emph{regularisation path}\index{regularisation path} of a penalised estimator is its solution, or its validation
score, as a function of the penalty $\lambda$; the penalty is chosen at the path's best purged cross-validated score.
\end{definition}

\begin{omfigure}
\begin{tikzpicture}
\begin{groupplot}[group style={group size=2 by 1,horizontal sep=1.6cm},width=6.6cm,height=5.2cm,
  tick label style={font=\footnotesize},label style={font=\small},grid=major,grid style={gray!25},table/col sep=comma,
  ymin=-0.02,ymax=0.42]
\nextgroupplot[xlabel={$\log_{10}\lambda$ (ridge)},ylabel={cross-validated $R^2$ (\%)},xtick={2,3,4,5},
  legend style={font=\footnotesize,at={(0.5,-0.33)},anchor=north,draw=none,fill=none},legend columns=2]
\addplot[omDef,very thick,mark=*] table[x=log10alpha,y=ridge]{figdata/ml/04-linear-and-regularised-baselines/path_ridge.csv};
\addplot[omProp,very thick,mark=triangle*] table[x=log10alpha,y=splines]{figdata/ml/04-linear-and-regularised-baselines/path_ridge.csv};
\legend{ridge,ridge on splines}
\nextgroupplot[xlabel={$\log_{10}\lambda$ (lasso)},xtick={-4,-3.5,-3,-2.5},
  xticklabel style={/pgf/number format/fixed,/pgf/number format/precision=1}]
\addplot[omThm,very thick,mark=square*] table[x=log10alpha,y=lasso]{figdata/ml/04-linear-and-regularised-baselines/path_lasso.csv};
\end{groupplot}
\end{tikzpicture}

\omcaption{\omterm{def:ml:linear-and-regularised-baselines:path}{Regularisation paths}: pooled out-of-fold R-squared of three purged folds of the twenty training years, as a
function of the penalty. Ridge peaks at $\lambda = 10^4$ (training folds of about 80\,000 rows), the lasso at $10^{-3}$; one step further the
lasso sets every coefficient to zero and scores zero. Data: \texttt{ml\_baselines.paths}.}\label{fig:ml:baselines:path}
\end{omfigure}

The paths are flat near their peaks and fall steeply on either side (\cref{fig:ml:baselines:path}). At its peak the
lasso keeps six coefficients; one step further it keeps none, and a model that predicts zero everywhere scores
zero, which is better than any of the unpenalised fits on the left of the figure. Least squares on 120\,000
rows and forty features looks like a well-determined problem; with the factor returns making the months, not the rows,
the effective sample (chapter~1), it is not, and the penalty that wins is large.

\begin{method}[Choosing a penalty on a panel]\label{met:ml:baselines:penalty}
\begin{enumerate}
\item Rank each feature by date; demean the target by date.
\item Build purged folds of whole dates (a date's rows never straddle folds), with an embargo of the label length.
\item Evaluate a logarithmic grid wide enough that both ends are worse than the middle; if the best value is at an end,
extend the grid.
\item Prefer the most penalised value within one standard error of the best (the path's plateau, not its peak).
\end{enumerate}
\end{method}

\section{Dimension reduction: principal-component and partial least squares}

\begin{definition}[Principal component regression, partial least squares]\label{def:ml:linear-and-regularised-baselines:pcr}
\emph{Principal component regression}\index{principal component regression} (PCR) regresses the target on the first
$k$ principal components of the features (Book~4, chapter~22). \emph{Partial least squares}\index{partial least squares}
(PLS) regresses it on $k$ components chosen for their covariance with the target: the first weight vector is
proportional to $X^\top y$, and each later one is found the same way after deflating $X$ by the components already
taken.
\end{definition}

\begin{proposition}[What each reduction keeps]\label{prop:ml:baselines:pls}
Among unit vectors $w$, $w\propto X^\top y$ maximises the sample covariance of $Xw$ with $y$, whereas the first principal
direction maximises the variance of $Xw$ whatever $y$ is. When the features are uncorrelated with equal variance, every
direction has the same variance and the principal components are arbitrary; PLS's first component is the least-squares
fit's direction up to scale.
\end{proposition}

\begin{proof}
$\widehat{\Cov}(Xw, y)\propto w^\top X^\top y \le \lVert w\rVert\,\lVert X^\top y\rVert$ by Cauchy--Schwarz, with equality
for $w\propto X^\top y$. If $X^\top X = cI$, the sample variance $w^\top X^\top Xw/n = c/n$ is the same for every unit $w$;
the least-squares coefficient is $(X^\top X)^{-1}X^\top y = X^\top y/c$.
\end{proof}

Cross-sectional ranks are close to that case: forty characteristics of equal variance and little correlation. PCR needs
ten components before it recovers the signal and still scores only 0.13\%; PLS's single component already does what
least squares does, and the two score alike (\cref{tab:ml:baselines:table}). On real characteristics, which are
correlated in blocks (value measures with each other, momentum measures with each other), PCR does better, and Gu,
Kelly and Xiu found PCR and PLS at 0.26\% and 0.27\%, above their elastic net's 0.11\%.

\section{Classification baselines}

\begin{definition}[Logistic regression]\label{def:ml:linear-and-regularised-baselines:logistic}
\emph{Logistic regression}\index{logistic regression} models $\P(y = 1\mid x) = 1/(1 + e^{-(b + x^\top\beta)})$ and fits
$(b,\beta)$ by maximum likelihood, usually with a ridge or lasso penalty on $\beta$.
\end{definition}

Labels that are signs (chapter~2) call for a classifier; its natural baseline is the \omterm{def:ml:linear-and-regularised-baselines:logistic}{logistic regression}. On the panel it
ranks stocks as well as the regressions do (IC 0.060, as for ridge) with a probability to spare: throwing away
the size of each return costs almost nothing when that size is mostly noise.

\begin{definition}[Basis expansion]\label{def:ml:linear-and-regularised-baselines:basis}
A \emph{basis expansion}\index{basis expansion} replaces each feature $x_j$ by functions $h_{j1}(x_j),\dots,h_{jm}(x_j)$
(splines, indicators of bins, polynomials) and fits a linear model on them; the model is additive in the features and
nonlinear in each.
\end{definition}

Ridge regression on quadratic splines with four knots per characteristic is the strongest linear baseline of the panel
(R-squared 0.35\%): it captures the square and the absolute value the truth contains, but not the interaction, which no
additive model can.

\begin{table}[htbp]
\centering\small
\begin{tabular}{@{}lcccccc@{}}
\toprule
model (parameter) & $R^2_{\mathrm{oos}}$ & IC & $t$ (HAC) & Sharpe gross & Sharpe net & turnover\\
\midrule
ordinary least squares & 0.25\% & 0.060 & 5.6 & 1.85 & 1.61 & 0.83\\
ridge ($\lambda = 10^4$) & 0.30\% & 0.060 & 5.6 & 1.83 & 1.59 & 0.83\\
lasso ($\lambda = 10^{-3}$) & 0.32\% & 0.066 & 6.7 & 2.17 & 1.91 & 0.80\\
elastic net & 0.32\% & 0.066 & 6.7 & 2.17 & 1.91 & 0.80\\
PCR (10 components) & 0.13\% & 0.043 & 4.1 & 1.23 & 0.99 & 0.77\\
PLS (1 component) & 0.25\% & 0.061 & 5.6 & 1.87 & 1.62 & 0.82\\
\omterm{def:ml:linear-and-regularised-baselines:logistic}{logistic regression} & 0.22\% & 0.060 & 5.7 & 1.84 & 1.59 & 0.84\\
ridge on splines & 0.35\% & 0.064 & 6.2 & 1.98 & 1.69 & 0.95\\
boosted trees & 0.47\% & 0.076 & 7.8 & 2.20 & 1.90 & 0.97\\
\midrule
ceiling & 0.71\% & & & & & \\
\bottomrule
\end{tabular}
\caption{Ten test years of the forty-characteristic panel with factor risk (500 stocks, twenty training years). Sharpe
ratios of the equal-weighted top-minus-bottom decile portfolio, annualised, net of 20 basis points per unit traded;
turnover is the share of the gross book traded each month. Data: \texttt{ml\_baselines.table}.}\label{tab:ml:baselines:table}
\end{table}

\section{Why the baseline often wins}

\Cref{tab:ml:baselines:table} ranks the models by R-squared, and the ranking by net Sharpe ratio is not the same. The
boosted trees forecast best (0.47\% against the lasso's 0.32\%), but their predictions change more from month to month
(turnover 0.97 against 0.80), their portfolio pays more in costs, and at the extremes of the ranking, where a decile
portfolio trades, the two models mostly agree. The lasso's net Sharpe ratio of 1.91 ties the trees' 1.90.

The comparison also depends on how much history there is (\cref{fig:ml:baselines:crossing}). Trained on the last two
years before the test, the trees' R-squared is $-0.71\%$, although their decile portfolio still earns a Sharpe ratio of
1.11 against ridge's 0.90: with little data the trees rank stocks tolerably and size their forecasts badly, which is the
difference between the IC and R-squared (\cref{prop:ml:why:ceiling}). They overtake ridge in R-squared between five and
ten years of history. On real data the argument has a second half: Kelly, Malamud and Zhou show that heavily
regularised models with more parameters than observations can predict better than small ones, so ``simple'' is not the
claim; ``regularised, and compared on equal terms'' is.

\begin{omfigure}
\begin{tikzpicture}
\begin{groupplot}[group style={group size=2 by 1,horizontal sep=1.6cm},width=6.6cm,height=5.2cm,
  xlabel={\small training window (months)},xtick={24,60,120,240},tick label style={font=\footnotesize},
  label style={font=\small},grid=major,grid style={gray!25},table/col sep=comma,xmin=10,xmax=254]
\nextgroupplot[ylabel={$R^2_{\mathrm{oos}}$ (\%)},ymin=-0.8,ymax=0.8,
  legend style={font=\footnotesize,at={(1.1,-0.33)},anchor=north,draw=none,fill=none},legend columns=3]
\addplot[omThm,very thick,mark=square*] table[x=months,y=ridge]{figdata/ml/04-linear-and-regularised-baselines/crossing.csv};
\addplot[omDef,very thick,mark=*] table[x=months,y=boosting]{figdata/ml/04-linear-and-regularised-baselines/crossing.csv};
\addplot[black!60,thick,dashed] table[x=months,y=ceiling]{figdata/ml/04-linear-and-regularised-baselines/crossing.csv};
\legend{ridge,boosted trees,ceiling}
\nextgroupplot[ylabel={net Sharpe ratio},ymin=0.6,ymax=2.4]
\addplot[omThm,very thick,mark=square*] table[x=months,y=ridge_sr]{figdata/ml/04-linear-and-regularised-baselines/crossing.csv};
\addplot[omDef,very thick,mark=*] table[x=months,y=boosting_sr]{figdata/ml/04-linear-and-regularised-baselines/crossing.csv};
\end{groupplot}
\end{tikzpicture}

\omcaption{Ridge and boosted trees trained on the last 2 to 20 years before the same ten test years. Left, R-squared
against zero; right, the net Sharpe ratio of the decile portfolio. The trees' ranking is useful early; their sizes are
not. Data: \texttt{ml\_baselines.crossing}.}\label{fig:ml:baselines:crossing}
\end{omfigure}

\section{Tutorial: the model to beat}\label{tut:ml:linear-and-regularised-baselines:tut}

\begin{tutorial}
\textbf{Goal.} Fit the linear baselines along their \omterm{def:ml:linear-and-regularised-baselines:path}{regularisation paths} through \texttt{firm.mlbase}, compare them with
boosted trees on one report, and vary the training window. \textbf{End state:} \cref{tab:ml:baselines:table},
\cref{fig:ml:baselines:path,fig:ml:baselines:crossing}.
\begin{tutsteps}
\item \textbf{The report}: R-squared against zero, the rank IC with its HAC $t$, and the decile portfolio gross and net.
\omcode{code/firm/mlbase/firm_mlbase.py}{61}{81}{The long-short decile portfolio and the standard report.}\label{lst:ml:baselines:report}
\item \textbf{The path}: purged folds of whole months, pooled out-of-fold R-squared.
\omcode{code/firm/mlbase/firm_mlbase.py}{84}{95}{Cross-validated regularisation path.}\label{lst:ml:baselines:path}
\item \textbf{Run} \texttt{ml\_baselines.paths()}, \texttt{table()}, \texttt{crossing()} and \texttt{fig\_baselines.py}.
\end{tutsteps}
\textbf{What to change next.} Set the factor volatility to zero and watch every Sharpe ratio double; add the
interaction $c_0c_3$ as a feature to ridge on splines and see how much of the trees' advantage it removes.
\end{tutorial}

\section{Build: the model harness}\label{bld:ml:linear-and-regularised-baselines:mlbase}

\begin{build}
\bfield{Purpose} Every cross-sectional model is compared with the same baselines, on the same splits, with the same
report.
\bfield{Interface} \texttt{rank\_features(X)}, \texttt{xs\_target(r, months)}, \texttt{fit(model, X, r, months,
target)}, \texttt{predict(model, X, months)}, \texttt{long\_short(pred, r, q)}, \texttt{report(pred, r, q, cost\_bp,
periods)} returning \texttt{r2}, \texttt{ic}, \texttt{ic\_t}, \texttt{ls\_sr}, \texttt{ls\_sr\_net},
\texttt{turnover}; \texttt{cv\_path(make, params, X, r, months, n\_folds, embargo)}; \texttt{Harness(make,
target).run(X, r, train, test)}.
\bfield{Rules} Any scikit-learn-style regressor plugs in; folds are whole dates; the report is computed on test dates
only; costs are charged on traded notional.
\bfield{Acceptance tests} \texttt{code/firm/mlbase/tests/}: ranks and demeaned targets; decile weights that sum to zero
with gross two; the harness finds a planted signal and reports gross above net; the path prefers a moderate penalty on
noisy data.
\bfield{Stretch} Value-weighted deciles; the report by size group; a random-feature ridge in the spirit of Kelly,
Malamud and Zhou as a further baseline.
\end{build}

\begin{omsources}
\item S.~Gu, B.~Kelly and D.~Xiu, ``Empirical asset pricing via machine learning'', \emph{Review of Financial Studies}
33(5), 2020.
\item I.~Welch and A.~Goyal, ``A comprehensive look at the empirical performance of equity premium prediction'',
\emph{Review of Financial Studies} 21(4), 2008.
\item B.~Kelly, S.~Malamud and K.~Zhou, ``The virtue of complexity in return prediction'', \emph{Journal of Finance}
79(1), 2024.
\item S.~de Jong, ``SIMPLS: an alternative approach to partial least squares regression'', \emph{Chemometrics and
Intelligent Laboratory Systems} 18, 1993.
\item A.~E.~Hoerl and R.~W.~Kennard, ``Ridge regression: biased estimation for nonorthogonal problems'',
\emph{Technometrics} 12(1), 1970.
\end{omsources}

\section{Exercises}

\begin{exercise}[$\star$]\label{exo:ml:linear-and-regularised-baselines:1}
A decile portfolio of 500 stocks holds 50 long and 50 short, equally weighted, gross 2. In a month 30 names change on
each side. What share of the gross book is traded, and what does it cost at 20 basis points per unit traded?
\end{exercise}

\begin{exercise}[$\star$]\label{exo:ml:linear-and-regularised-baselines:2}
A long-short portfolio returns 1.6\% a month on average with a monthly volatility of 2.5\%. What is its annualised
Sharpe ratio? And after a monthly cost of 0.2\%?
\end{exercise}

\begin{exercise}[$\star$]\label{exo:ml:linear-and-regularised-baselines:3}
With two features of equal variance whose correlation is 0.9, which direction does PCA's first component take, and
when is it the right direction for a regression?
\end{exercise}

\begin{exercise}[$\star\star$]\label{exo:ml:linear-and-regularised-baselines:4}
Why do the boosted trees and the lasso earn the same net Sharpe ratio although the trees' R-squared is half as large
again?
\end{exercise}

\begin{exercise}[$\star\star$]\label{exo:ml:linear-and-regularised-baselines:5}
Trained on two years, the trees' R-squared is $-0.71\%$ and their decile Sharpe ratio 1.11. Reconcile the two numbers
with \cref{prop:ml:why:ceiling}.
\end{exercise}

\begin{exercise}[$\star\star$]\label{exo:ml:linear-and-regularised-baselines:6}
\emph{Find the flaw.} ``Our deep model's R-squared beats the published linear benchmark by 0.2 points, so it is better.''
The published benchmark was run on another period and universe.
\end{exercise}

\begin{exercise}[$\star\star\star$]\label{exo:ml:linear-and-regularised-baselines:7}
\emph{Coding.} Rerun \texttt{ml\_baselines.table} with the factor volatility set to zero (\texttt{STYLE = 0}). Report
ridge's and the trees' net Sharpe ratios and explain why they change so much while R-squared hardly does.
\end{exercise}

\begin{exercise}[$\star\star\star$]\label{exo:ml:linear-and-regularised-baselines:8}
Show that for an orthonormal design the lasso sets to zero every coefficient whose least-squares value is below
$\lambda$ in absolute value, and use it to explain the lasso path's drop to zero in \cref{fig:ml:baselines:path}.
\end{exercise}

\section{Problem: The Model to Beat}

\begin{problem}[{Weekend problem --- a quarter's work against a second's}]\label{pb:ml:linear-and-regularised-baselines:1}
The chapter's panel: 500 stocks, forty characteristics, factor risk, twenty training years and ten test years.

\medskip\noindent\textbf{Part I --- The panel.}
\begin{enumerate}
\item What is the ceiling, and which characteristics carry signal?
\item Why were factor returns added, and what do they do to Sharpe ratios?
\item What equal-weighted decile Sharpe ratios did Gu, Kelly and Xiu report?
\item Why does least squares on 120\,000 rows need a penalty?
\end{enumerate}

\medskip\noindent\textbf{Part II --- The linear family.}
\begin{enumerate}[resume]
\item Which penalties do ridge and the lasso choose, and what does the lasso score one step further?
\item Why does PCR need ten components and still score 0.13\%?
\item Why does PLS with one component score like least squares?
\item What does the spline expansion capture, and what can it not?
\end{enumerate}

\medskip\noindent\textbf{Part III --- Trees against the baseline.}
\begin{enumerate}[resume]
\item Give the trees' and the lasso's R-squared, IC and net Sharpe ratio.
\item Why is the ranking by Sharpe ratio different from the ranking by R-squared?
\item How long a history do the trees need to overtake ridge in R-squared?
\item What do the trees earn with two years of history, and why?
\end{enumerate}

\medskip\noindent\textbf{Part IV --- The verdict.}
\begin{enumerate}[resume]
\item State the \emph{named result}: ridge's R-squared and IC against the trees', and the history at which the trees
overtake.
\item Which model would you run in production while the trees are validated, and why?
\item What further baseline would the virtue-of-complexity result suggest?
\item How many months of live results would tell the lasso from the trees in net Sharpe ratio?
\item What would make PCR the right baseline?
\item Which report column would you show a portfolio manager first?
\item What is lost by converting returns into signs?
\item In one sentence: what is a baseline for?
\end{enumerate}
\end{problem}

\section{Interview questions}

\begin{interviewq}[$\star$ \iqroles{researcher, mle}]\label{iq:ml:linear-and-regularised-baselines:1}
You are given a new cross-sectional return dataset. What is the first model you fit, and how do you set it up?
\end{interviewq}

\begin{interviewq}[$\star$ \iqroles{researcher}]\label{iq:ml:linear-and-regularised-baselines:2}
Ridge or lasso for forty ranked stock characteristics? Why?
\end{interviewq}

\begin{interviewq}[$\star\star$ \iqroles{researcher}]\label{iq:ml:linear-and-regularised-baselines:3}
What is the difference between \omterm{def:ml:linear-and-regularised-baselines:pcr}{principal component regression} and \omterm{def:ml:linear-and-regularised-baselines:pcr}{partial least squares}, and when does it matter?
\end{interviewq}

\begin{interviewq}[$\star\star$ \iqroles{researcher, trader}]\label{iq:ml:linear-and-regularised-baselines:4}
A model improves R-squared by 50\% but not the portfolio's net Sharpe ratio. How can that be?
\end{interviewq}

\begin{interviewq}[$\star\star$ \iqroles{mle}]\label{iq:ml:linear-and-regularised-baselines:5}
How do you choose the penalty of a ridge regression on a panel of stocks, concretely?
\end{interviewq}

\begin{interviewq}[$\star\star\star$ \iqroles{researcher}]\label{iq:ml:linear-and-regularised-baselines:6}
Derive the first partial-least-squares direction, and say why it can beat the first principal component.
\end{interviewq}
